% ----------------------------------------------------------
% Fundamentação Teórica
% ----------------------------------------------------------
\chapter{Revisão da Literatura}

O aumento do número de pessoas que vivem sozinhas tem sido descrito como um fenômeno global com impactos relevantes na saúde e no bem-estar. Diversos estudos apontam que essa condição está associada a maior risco de hospitalizações de emergência, especialmente entre idosos, devido à ausência de suporte imediato em situações críticas \cite{kandler2017}.

No Brasil, pesquisas publicadas na Revista Brasileira de Enfermagem evidenciam que idosos que vivem sozinhos apresentam maior vulnerabilidade social e limitações relacionadas ao suporte familiar \cite{perseguino2017}. Além disso, estudos internacionais apontam maior incidência de quedas e outros eventos adversos entre idosos residentes em moradias unipessoais \cite{nakamura2017}. Esses dados reforçam a necessidade de soluções voltadas ao monitoramento e suporte dessa população.

Nesse contexto, tecnologias como casas inteligentes e sistemas baseados em Internet das Coisas (IoT) têm ganhado destaque, permitindo o acompanhamento remoto de sinais vitais e padrões de atividade \cite{smarthomehealth2019,predictivehealthcare2020}. Embora promissoras, tais soluções ainda enfrentem desafios relacionados à segurança da informação, ao custo e à acessibilidade, o que abre espaço para iniciativas como o Linha Vital, que busca oferecer uma alternativa prática, não invasiva e centrada no uso doméstico.

\section{Histórico da Tecnologia no Monitoramento da Saúde}

O avanço das tecnologias aplicadas à saúde tem possibilitado o desenvolvimento de soluções voltadas ao monitoramento remoto e à assistência de pessoas em situação de vulnerabilidade, especialmente idosos e pessoas com deficiência. Nesse contexto, as chamadas casas inteligentes e sistemas de monitoramento domiciliar surgem como alternativas tecnológicas para promover maior segurança e qualidade de vida no ambiente doméstico \cite{smartliving2018,ambientassistedliving2016}.

Estudos indicam que essas tecnologias ganharam relevância à medida que se buscava melhorar o bem-estar de pessoas com condições crônicas ou limitações funcionais, utilizando sensores e sistemas automatizados para acompanhamento contínuo. Soluções baseadas em monitoramento inteligente passaram a ser analisadas de forma sistemática, buscando compreender seu impacto na saúde, na segurança e na autonomia dos usuários \cite{ambientassistedliving2016}.

Apesar dos avanços, ainda há uma quantidade limitada de estudos que mensuram de forma concreta os benefícios dessas soluções em larga escala. Dessa forma, observa-se que o uso de tecnologias para assistência em saúde evoluiu de abordagens tradicionais para soluções inteligentes baseadas em automação residencial e monitoramento remoto, consolidando-se como uma área de crescente interesse acadêmico e tecnológico.

Nesse cenário, o projeto Linha Vital posiciona-se como uma alternativa prática e acessível, alinhada às tendências de inovação, mas voltada especificamente ao uso doméstico e cotidiano.

\section{Atualidade da Tecnologia no Monitoramento da Saúde}

Atualmente, o uso de tecnologias para monitoramento remoto de saúde tem se intensificado, especialmente com a incorporação da Internet das Coisas (IoT), que permite a coleta e análise de dados em tempo real. Essas soluções possibilitam o acompanhamento contínuo de usuários fora do ambiente hospitalar, contribuindo para uma resposta mais rápida a situações de risco \cite{iothealthcare2020,predictivehealthcare2020}.

Pesquisas apontam que sistemas baseados em IoT são capazes de monitorar sinais vitais e atividades dos usuários, enviando alertas para familiares ou profissionais de saúde em caso de anormalidades, reduzindo a necessidade de deslocamentos até unidades médicas \cite{iothealthcare2020}. Além disso, estudos evidenciam que essas tecnologias possuem potencial para melhorar a qualidade de vida dos usuários ao proporcionar maior independência e segurança no ambiente doméstico \cite{smarthomehealth2019}.

Entretanto, apesar dos avanços, ainda existem desafios importantes relacionados à segurança da informação, à confiabilidade dos sistemas e à necessidade de maior maturidade tecnológica para aplicação em larga escala \cite{predictivehealthcare2020}. Nesse contexto, o projeto Linha Vital busca oferecer uma alternativa prática e acessível, centrada no uso doméstico, que alia os benefícios da IoT à simplicidade necessária para atender indivíduos em situação de vulnerabilidade.

\section{Outros Contextos da Tecnologia no Monitoramento da Saúde}

Além do uso no contexto de saúde domiciliar, tecnologias de monitoramento e automação vêm sendo aplicadas em diferentes cenários, como acessibilidade, inclusão social e assistência a pessoas com deficiência. Pesquisas demonstram que sistemas inteligentes têm sido desenvolvidos para monitorar ambientes e detectar situações de risco, utilizando sensores e dispositivos conectados para coleta de dados e emissão de alertas, ampliando a segurança de indivíduos em situação de vulnerabilidade \cite{reboucas2019}.

Da mesma forma, soluções baseadas em automação residencial têm possibilitado que pessoas com limitações físicas ou sensoriais conquistem maior autonomia no cotidiano por meio do controle de dispositivos e do acompanhamento de atividades, promovendo inclusão e independência funcional \cite{automacaoiot2019}.

Entretanto, estudos apontam que muitas dessas iniciativas ainda apresentam barreiras significativas, como custos elevados de implantação e complexidade técnica, fatores que dificultam sua adoção em larga escala e restringem o acesso de parte da população \cite{reboucas2019}.

Nesse cenário, observa-se que a literatura evidencia tanto o potencial transformador dessas tecnologias quanto os desafios que precisam ser superados para que se tornem soluções efetivamente inclusivas. É nesse ponto que o projeto Linha Vital se insere, ao propor uma alternativa prática e acessível, que alia os benefícios do monitoramento inteligente à simplicidade necessária para atender indivíduos em situação de vulnerabilidade, sem exigir infraestrutura complexa ou investimentos elevados.